\section{Integral level moves and extremal evaluation on the given geometric half}

% BEGIN G77D1
\paragraph{G77D1. The integral finite-group lattice on the original record.}
Let $B$ be any original $\mathrm{GeoHalf}(T,R,\mathrm{ph})$ and
write $w=B.\mathrm{fw}$. Then $R.\mathrm{LIntegral}$ holds.
This conclusion already follows from $w$, the base level-action
pins and $B.\mathrm{rhoGeo}$.

For row $i$, let $C_i$ be the lower algebraic cusp space over
$K=T.Q.c.F$. The tie supplies an isomorphism $e_i:C_i\to W_{K,i}$
onto the original rational coordinate space, carrying the
all-valuations condition onto membership in $M_{K,i}$.
Take a move $g$ of the marked level. In the branch
$\mathrm{ActKeeps}$ the section $g f$ is again a cusp section and
the coordinate move on $e_i(f)$ is $e_i(g f)$. Consider one actual
valuation test for $g f$. Its point $y$ over the fraction field
comes from a test ring, and its frames come from the integral
model over the given valuation ring. Apply the same marking move
on the test ring and then extend to the fraction field. Naturality
of the level action identifies the resulting point with $g y$.
The integral group scheme, polarization and valuation ring can be
retained. Compose its generic-fiber comparison with the prescribed
change of marking. This carries each Lie basis to the basis specified
at $g y$. The isomorphism of rigidified theta lines carries its
generator to the new generator. Indeed, on a trivializing open the
new section is the old section times the transition unit of the
isomorphism, and these local expressions glue by its cocycle law.
Thus $g y$ again has an integral frame over the same ring. Evaluation
by pullback gives $(g f)(y)= f(g y)$, so the normalized valuation
is nonnegative whenever $f$ passes all tests. This works separately
at every test, including evaluations equal to zero. Both rational
space and integral lattice are therefore carried into themselves.
In the other branch $B.\mathrm{rhoGeo}$ specifies the identity
matrix, which preserves both. This proves the full two-part
$\mathrm{MKStable}$ at $i$. Combine it with $w.\mathrm{ks}(i)$
and the full lattice of M76L2 to obtain the three clauses of
$R.\mathrm{LIntegral}$ for every row.
% END G77D1

% BEGIN G77D2
\paragraph{G77D2. Evaluation lifts at every good embedding.}
Keep the geometric half, assume the level law
$R.\mathrm{toLevCh.Is}(A.\mu_0)$ and $R.\mathrm{IsModels}$, and assume $T.\mathrm{pb}=R.C.P.\mathrm{pb}$. For
$i$ and $\sigma:L\hookrightarrow\mathbb C$ put
$\alpha=\sigma|_K$. At every good $\alpha$ with the prescribed
identification of Hermitian spaces there is a unique
$\alpha$-semilinear map $j_\sigma:C_i\to H_{i,\sigma}$
with $\mathrm{ev}_{v_0}j_\sigma(f)=\phi_C(f)$.
This is the field $\mathrm{JEvExists}$.

Here the finite level in the target model is the principal kernel.
The two marked primes agree in order, both self-dual frames realize
the completed reference lattice, and reducing the frame modulo the
level tests exactly membership in this kernel. On the trivial
tuple its level character is one. At a good embedding the parameter
tie puts the positive signs in the holomorphic order and the
weight tie gives the corresponding integral extremal weight.
The core laws identify the component Gram form with the prescribed
one, match the finite level and send adelic points to frames in the
period domain. Write $P$ for a generating tuple of bounded
holomorphic cusp components in the definition of $\mathrm{holSub}$.
Its adelic function is
\[
 A(P)(h)=c_{a(h)}\mu(h)P_{a(h)}(x(h)),\qquad
 (a(h),x(h))=\mathrm{dict}(h).
\]
The rational and finite-level transformation laws follow by changing
the dictionary mover: the arithmetic stabilizer has trivial weight
character and $\mu$ is trivial on the principal kernel. The
component bounds give a single bound since there are finitely many
components and the multipliers have norm one. For the holomorphic
law above $h$, use its fixed rational mover and let the real unitary
matrix act on the orthonormal reference frame. The corresponding
frame variable runs through the Gram domain. Multiply the component
section there by the determinant powers from the weight tie,
normalizing by the same powers at the reference frame. These
determinants never vanish there. The triangular transformation law
is the extremal Borel weight and the unitary determinant factor
is precisely $\mu$ on the real component. This gives the required
holomorphic local representative and continuity; the finite level
is open.

For a rational isotropic direction $u$, transport it by the rational
mover and multiply by a positive integer $N$ clearing the
denominators to get the chart direction. Conjugation takes the
transvection to the chart flow, scaling its real parameters at
$\tau$ by the chart scalar times $N^2$. All these factors are
nonzero. Transvections have determinant one, so the character
factor is constant on this flow. The chart cusp law gives a full
box of real periods whose mean is zero at every starting frame.
To pass to the adelic constant term, average this identity over
the rational-transvection quotient with its finite Haar measure.
Translation invariance gives the same quotient integral at each
parameter. Fubini is valid by boundedness and finite measure, and
the period box has positive volume. The quotient integral vanishes.
Thus $A(P)$ satisfies $\mathrm{HolForm}$ and $\mathrm{ZCusp}$
on the target level.

FV8 supplies a finite pair. By SPR21 there is a smooth cuspidal
multiplicity map $U$ with $U(v_0)=A(P)$. The model identifies the
range of its map $e$ with the span of exactly these multiplicity
maps, so $U=e(x)$ for some $x\in H_{i,\sigma}$. The range of
evaluation is a complex subspace. It therefore contains $A$ of
the entire component span, including the reading of each algebraic
cusp section, and still contains it after multiplication by the
prescribed Lie factor. The latter is exactly $\phi_C$. The
unit-vector norm argument of T76R1 proves injectivity from the
model inner product and $w$ alone. Choose the unique preimage of
$\phi_C(f)$ for each $f$. Applying evaluation to sums and
$K$-multiples proves additivity and $\alpha$-semilinearity.
% END G77D2

% BEGIN G77D3
\paragraph{G77D3. Conjugate evaluation with one unit tuple.}
Under the hypotheses of G77D2, let $\sigma$ have nongood
restriction and put $\sigma'=\overline\sigma$. Its restriction
is good. There is an antilinear bijection
$J:H_{i,\sigma'}\to H_{i,\sigma}$ and one smooth unit tuple $v$
at $\sigma$, chosen before $f$, such that
\[
 \mathrm{ev}_{v}(Jy)=\overline{\mathrm{ev}_{v_0(i,\sigma')}(y)}
 \quad(y\in H_{i,\sigma'}).
\]
In particular $Jj_{\sigma'}$ is $\sigma|_K$-semilinear with the
specified conjugate reading. This proves $\mathrm{JEvBar}$.

The parameter at $\sigma'$ is the dual parameter at $\sigma$.
Conjugate the local unitary models at $\sigma$: complex
conjugation reverses the compact weights, interchanges raising
and lowering equations, and carries the smooth extremal vector
to an extremal vector for the dual parameter. The resulting
holomorphic model and the stored one at $\sigma'$ have the same
extremal coefficient. Indeed HCU15 and HCU7 give proportionality
from the holomorphic Borel law, and value one at the identity
fixes the constant for their unit vectors. On finite sums of
translates of the extremal vector the two Gram matrices are
therefore equal. Their spans are dense by irreducibility; completion
gives a unitary intertwiner. Composing these identifications
gives local antiunitary intertwiners $c_\tau$ from the stored
model at $\sigma$ to that at $\sigma'$. Take
$v_\tau=c_\tau^{-1}v_0(i,\sigma')_\tau$. Intertwining and
real smoothness preserve smooth vectors, and isometry gives norm one.

For a multiplicity map $U$ at $\sigma'$, the rule
$u\mapsto\overline{U((c_\tau u_\tau)_\tau)}$ is complex
multilinear. Conjugation commutes with right translation and
real differentiation, preserves growth and smooth-vector bounds,
and conjugates the cusp integrals to zero. The principal-level
character is one. Thus conjugating in this fashion sends the
span of multiplicity maps to the corresponding span at $\sigma$.
The inverse intertwiners give the opposite inclusion.
Using the range clauses of the stored models now gives a preimage
at $v$ of each conjugate evaluation and conversely. Evaluation is
injective at either tuple by the same norm argument, which requires
only unit components. The two choices of preimages are inverse
since complex conjugation squares to the identity. Their addition
and scalar laws follow after evaluation. This constructs the
asserted $J$ on the original spaces.
% END G77D3

% BEGIN G77D4
\paragraph{G77D4. Component constants and the level action at all embeddings.}
Under the hypotheses of G77D2, $B$ provides the original
$\mathrm{GeoAct}$. We construct the constants individually at
each good embedding $\alpha$ of $K$; their transformation law
uses $\alpha(y_g)$, where $y_g\in K$ has distinguished value
$\chi_C(g)$, the product of the marked determinant characters.

A lift $k$ of $g$ acts on the finite components on the right;
write $a\mapsto ag$. This is independent of lift: two lifts
differ by the principal kernel. The period markings for $h,hk$
are related by the rational isometry obtained by composing their
dictionary movers, since $k$ preserves the finite reference lattice,
induces the prescribed marked action there, and has trivial
infinite part. In particular the geometric change of marking is
the one induced by $k$. On the algebraic sections it gives the
pullback $f\mapsto gf$. The section laws are natural under this
change. Its chart reading is carried to the new component by
the rational isometry, up to the change of the normalized
theta generator. The latter is a unit in absolute value on
orthonormal frames: the two pulled rigidified lines and their
polarizations agree, and integration of the squared norm of their
generators against the normalized Gaussian metric gives the same
value. This comparison applies to all sections with the given
exponent. Removing a triangular factor gives the same conclusion
on general frames since its Borel multiplier occurs on both sides.

Write $r_f(h)$ for the raw frame reading. For fixed $a,g$ the
ratio just found can be chosen as a single number $z(a,g)$
independently of $h,k,f$, with
\[
 |z(a,g)|=1,\qquad r_{gf}(h)=z(a,g)r_f(hk).
\]
Indeed two rational movers inducing the same change of marking
differ by an arithmetic stabilizer with trivial weight character.
The remaining ratio depends initially on the frame. If one
denominator is nonzero, the ratio of the two holomorphic readings
there has constant modulus, hence is constant on an open
neighborhood. The identity theorem on the connected frame domain
extends the simultaneous equality for every section to the
component. If all readings vanish take $z=1$. This description
also preserves holomorphy, the weighted bound and the vanishing
of constant terms: the isometry transports period boxes for each
rational isotropic direction to period domains for its image.
It follows at every good embedding that $gf$ is a cusp section.
Composition of pullbacks therefore gives a representation on $C_i$.

For a move fixing its component the two raw readings are equal,
including the phase. Here is why the unit in the preceding
calculation can be set equal to one on such a move. The rational
mover is a unitary automorphism of the period lattice. Its
determinant is integral together with its inverse, has modulus
one in every complex embedding, and hence is a root of unity.
The host has $4\mid[E:\mathbb Q]$. The squared theta transformation
of NFL2--NFL7 compares the two normalized generators on this
same component. In this application the rigidified unit is
the ratio of the pullback and original trivializations of the
theta line; its value at the origin is one. The two normalizing
constants squared are the same fixed nonzero multiple of their
respective period determinants. For the sign in this
transformation, write $K_0=\mathbb Q(\sqrt{-d})$ for the host
quadratic field. The relative norm to $K_0$ of the determinant
root of unity equals one. Indeed its odd part has norm of odd
order in $\mu(K_0)=\{\pm1\}$. The $2$-power cyclotomic fields
are disjoint from $K_0$ since $d$ is an odd prime ramifying in
$K_0$. The norm of a primitive $2^s$-th root for $s\ge2$ in the
cyclotomic field equals one, and the norm of $-1$ equals one
since $[E:K_0]$ is even. In the adapted symplectic basis the
host idempotent $e=w_0^4$ modulo eight projects onto the first
block. The mover commutes with $e$, so its off-diagonal blocks
vanish modulo eight. The determinant of the first block is the
determinant of $e\delta+(1-e)I$ modulo eight. Computing it by
restriction of scalars first to $K_0$ gives one modulo eight:
on the $e$ summand its determinant over $K_0$ is the relative
norm just computed, and on $1-e$ it is one. These congruences
give the level-eight sign and the determinant cancellation in
the squared theta transformation.
Thus the relative generator scalar has square one. Its contribution
at the even exponent $n$ equals one. The Lie frames are the
transports by the same isometry. This proves equality at the
stabilizer, independently of zeros of any one cusp section.

The residue calculation of X76R1 gives a finite idele determinant
root with trivial infinite part and distinguished character
$\chi_C(g)^{-1}$. On this same idele the core law
$\mathrm{muConj}$ compares the two good embeddings by the value of
an element of $K$. Inverting it yields $y_g$ with distinguished
value $\chi_C(g)$ and $\alpha(y_g)\mu_\alpha(k)=1$.
Injectivity at the distinguished embedding gives $y_{gg'}=y_gy_{g'}$.
Consequently it suffices to arrange $c_a r_{gf}(h)=c_{ag}r_f(hk)$.

Call a component active when some reading of a cusp section at an
adelic point in it is nonzero. Activity is preserved in both
directions by the displayed equation and the inverse level move.
On each active orbit compare the equation for $gg'$ with the two
successive equations for $g,g'$, choosing an inverse-translated
section so that the common factor is nonzero. Cancellation gives
$z(a,gg')=z(a,g)z(ag,g')$. The stabilizer calculation gives
$z(a,g)=1$ when $ag=a$. Set $c_{ag}=z(a,g)$ from a fixed
representative $a$. This is well defined and gives the required
identity on that orbit. Set $c=1$ on inactive orbits. Multiplying
by $\mu_\alpha(h)$ now proves $\mathrm{EqAt}$ with these unit
constants. The implication $B.\mathrm{equivAll}$ installs the
law for the \emph{prescribed} $T.c_E$ at this same $\alpha$.

To check $\mathrm{IsGeoAct}$ use the representation $f\mapsto
gf$ just constructed. At a good restriction the $j_{\rm Top}$
evaluation is G77D2, with the prescribed Lie factor independent of
$h$. Evaluate the level action using its lift $k$; the target
multiplicity map acts by $h\mapsto hk$. The preceding identity
then gives the required scalar $\sigma(\chi_{\rm Top}(g))$,
since the field character defined by $y_g$ has the required
distinguished values. At a nongood restriction evaluate with
the unit tuple of G77D3 and conjugate the same equality at
$\overline\sigma$. If the defining space identification is absent,
the top map is zero by definition. In the nonzero branches equality
of the evaluations implies equality of the multiplicity maps:
the difference is in their span and its inner product at a unit
tuple detects its norm, as in G77D2. Finally the integral part
of $\mathrm{IsGeoAct}$ is the integral-frame calculation in
G77D1 applied to this representation. This proves both fields.
% END G77D4

% BEGIN G77D5
\paragraph{G77D5. Why the short theta generators in these weights are cusp sections.}
Fix a row and suppose $8D.R.f\,D.R.r+4\le D.R.nV$.
The identity-cover short sections of B76S2 in row weight
$(a,n)=(\mathrm{wtA}(\kappa_i),\kappa_i.n)$ belong to the original
lower cusp space $C_i$. In particular this is an assertion about
their readings at \emph{every} good embedding.

First put $p=D.R.k+i$ for this row. At any infinite place the
weight parameters in decreasing order satisfy
$\lambda_s\ge\lambda_{p-1}+p-1-s$ for $s<p$, and
$\lambda_{p-1}=-b_i$. Thus the positive block weight at $s$,
the nonnegative integer part of
$\lambda_s+e_i-(p-1-s)-(D.R.t+1)/2$, is at least one by
$e_i-b_i\ge(D.R.t+3)/2$. In particular its last entry is
positive. The dictionary carries this last-entry condition to
one block of the row-highest harmonic polynomial.

Fix now any good embedding and any of its components. S76G2--S76G4
identify the short-rule reading there with the corresponding
harmonic theta series, up to a nonzero constant. For the present
argument one may use the component identity alone: choose a
symplectic basis for its unimodular period lattice, expand the
finite Schwartz function in point characteristics, and differentiate
the resulting classical theta functions with the specified harmonic
polynomial. The result is a finite sum. Each term, including its
Gaussian correction and period determinant, has the prescribed
Borel weight law. The arithmetic group law for the sum also follows
directly from the algebraic section.

Here are the two estimates needed for cuspidality. Write
$\ell(x)$ for the first minimum of the positive Gram majorant at
$x$, in length units. On an orthonormal frame the weighted norm of
each undifferentiated, metrically normalized theta generator has
bounded mean square for torus probability measure. The Haar
covolume in the orthonormal coordinates is fixed by the principal
polarization. On a ball of radius proportional to
$\min(1,\ell(x))$ the covering onto the torus is injective.
In a trivialization centered at its center the Gaussian metric
and its inverse are bounded there by constants depending on the
dimension. The submean inequality therefore bounds the local
supremum on its half-ball by a fixed inverse power of that radius.
Cauchy's inequalities on a smaller ball do so for any of the
finitely many derivatives. Apply this on the product torus for
the $n$ generators; its shortest lattice length is the same.
This gives $C(1+\ell(x)^{-1})^A$ for the weighted norm of the
sum. The constants may depend on the component and section.
Gram normalization by a triangular change of frame preserves
the weighted norm, hence gives this bound on the whole domain.

For a rational isotropic direction $u$, take the primitive
saturation of its period sublattice and extend it to a symplectic
basis. This is possible by unimodularity, dividing out each
successive primitive symplectic pair. The theta transformation under
the basis change replaces the characteristics by transformed ones
and a nonzero scalar. In the adapted basis the period matrix along
the transvection changes by a sum of rank-one quadratic forms in
its real parameters $s_\tau$. Thus a summand indexed by the
shifted integral rows $m_b+a_b$ has phase
$\exp(4\pi i\sum_\tau \nu_\tau s_\tau)$, where
\[
 \nu_\tau=\frac{c_\tau^2}{4}
       \sum_{b=1}^n |L_{u,\tau}(m_b+a_b)|^2 .
\]
Here $L_{u,\tau}$ is the linear functional obtained by pairing
the row with the period of $u$ in the dual half of the basis,
and $c_\tau$ is the chart's real scalar. Positive imaginary
period matrix gives absolute convergence after any fixed number
of differentiations, by Gaussian decay times a polynomial.

If the frequency vector is zero, every displayed linear pairing
vanishes. To see the effect on the polynomial differential operator,
make a unipotent triangular frame change eliminating the complementary
component of the period of $u$: in the two blocks the isotropy
equation is the compatibility condition for the off-diagonal
matrix $K$ in that triangular change. In block notation the two
equations are $K w_2=w_1$ and $l_1^{\mathsf T}K=-l_2^{\mathsf T}$.
Here $w_2,l_1$ are nonzero and $l_1^{\mathsf T}w_1
=-l_2^{\mathsf T}w_2$. Define $K$ on the line of $w_2$ by the
first equation, and extend on any complement by the second,
using a vector on which $l_1$ takes value one. At the chosen place the projection to each
nonempty definite block of the isotropic direction is nonzero
(otherwise definiteness and isotropy would force the whole vector
to vanish). In this gauge the Gaussian quadratic part and the
linear Fourier part are invariant under translation in the
projected period direction in each column of the relevant block.
For the quadratic part this follows from the eliminated complementary
component, and for the linear part from the vanishing pairings.
Consequently the matrix of row differentiation directions has a
nontrivial linear combination annihilating the summand, column by
column. The row-highest polynomial is a linear combination of
products of flag minors: apply the highest-weight polynomial
realization and expand changes of columns by Cauchy--Binet.
The positive last entry puts a full-block minor in each product.
Its determinant of commuting partial derivatives vanishes after
replacing one row by the annihilating combination. Thus the
coefficient of the zero frequency vanishes. A definite component
with no isotropic direction has no such cusp test.

For clarity, polynomial growth alone would give too weak a norm
bound here. In a cusp the same nonzero-frequency expansion gives a
uniform bound as follows. Put $L=\ell(x)^2$. For small $L$ choose
a nonzero integral $u$ of squared length less than $2L$. It is
isotropic: integrality of the lattice Hermitian pairing up to a
fixed denominator and the embeddingwise majorant inequality would
otherwise give a nonzero algebraic norm bounded above by
a constant times a positive power of $L$. In the transvection
coordinate put $t_\tau=(c_\tau\operatorname{maj}_\tau(x,u))^{-1}$.
Then $t_\tau\ge (2L)^{-1}$. For this fixed chart and weight let
$N=X.m\,X.\mathrm{chiOrd}(\lambda)$ and $E^+$ be the real subfield.
The period lattice is the real embedding of $N\mathcal O_{E^+}$:
these are exactly the parameters generated by the real algebraic
integers in the definition of $X.\mathrm{periods}(\lambda)$.
It is independent of $x$ and $u$. For the phases
$\exp(4\pi i\langle\nu,s\rangle)$ a nonzero Fourier coefficient
therefore has $2\nu$ in the dual lattice. The diagonal period $N$
gives $2N\sum_\tau\nu_\tau\in\mathbb Z$. Nonnegativity of each
$\nu_\tau$ and the missing zero coefficient give
$\sum_\tau\nu_\tau\ge1/(2N)$ on the support. One can obtain these
support statements even if individual characteristics are not
periodic at level $N$: regroup their absolutely summable real-flow
series by frequency. Translation by any period leaves the function
unchanged, so uniqueness of absolutely convergent exponential
expansions kills every nondual frequency. This uniqueness follows by
averaging against the opposite exponential on expanding cubes and
using dominated convergence.

For this convention the coefficient integral contains the factor
$\exp(-4\pi i\langle\nu,s\rangle)$. Under
$s\mapsto s-i t/2$ its modulus acquires the factor
$\exp(-2\pi\sum_\tau\nu_\tau t_\tau)$. To justify the simultaneous
shift use the straight homotopy $s-i b t/2$ for $0\le b\le1$.
The full segment lies in the tube
$\operatorname{Im}s_\tau>-t_\tau$ for every $\tau$. For each
coefficient the corresponding holomorphic top-degree form is
periodic in the real period lattice. Stokes' theorem on the
period parallelepiped crossed with this segment cancels opposite
vertical faces, giving equality of the two integrals. This uses
the actual possibly oblique period parallelepiped.

Fix $B$ bounding all coordinate absolute values on this one
parallelepiped. The moderate bound on the translated frames is at
most $C'(1+\ell(x)^{-1})^{A'}$ in the original normalization,
with constants independent of $u,x$. Indeed in a Gram-orthonormal
basis the rank-one nilpotent term for real translation at $\tau$
has norm bounded by a fixed multiple of $1/t_\tau$. Both the
translation and its inverse are bounded by $1+B/t_\tau$ in the
majorant metric (with the chart scalar absorbed in $t_\tau$).
On the hyperbolic plane of $u$ the imaginary shift changes the
positive pivot by $q=1+\operatorname{Im}s_\tau/t_\tau$.
In complementary isotropic directions the majorant changes by
$q$ and $q^{-1}$, and is unchanged on their orthogonal complement.
Here $1/2\le q\le1$. Thus the majorant quadratic forms are
uniformly comparable, for example within factors
$2(1+B)^2$ when all $t_\tau\ge1$. Taking minima over the
same lattice gives uniform comparability of squared first minima.
The fixed weight normalization is a product of fixed powers
of Gram minors, so the bounds for their ratios are uniform as
well. Nothing in the constants requires choosing finitely many
directions $u$.
Summing the coefficient bounds gives, for small $L$,
\[
 C''(1+L^{-1/2})^{A'}
     e^{-\pi/(4NL)}
        \sum_{\nu\ge0\ \mathrm{in\ the\ period\ dual}}
                      e^{-\pi\sum_\tau\nu_\tau}.
\]
The last series converges by lattice counting. Exponential decay
bounds the first factor uniformly, and for $L$ bounded below the
moderate estimate itself suffices. This proves the weighted bound
and zero constant term together, hence membership in the required
component cusp space. The embedding and component were arbitrary,
so the algebraic sections in question lie in $C_i$.
% END G77D5

% BEGIN G77D6
\paragraph{G77D6. Spanning the actual multiplicity space by algebraic readings.}
Assume the hypotheses of G77D2 and the rank inequality of G77D5.
Then $\mathrm{JUnif}(T,R,\kappa,v_0)$ holds. The analytic input
in the direction from multiplicity vectors to component functions
is a bound for each individual smooth cusp form. RPI11 supplies
this bound on the original model.

Fix a good restriction $\alpha$ and $y\in H_{i,\sigma}$.
Its image under $e$ belongs to the span of smooth multiplicity
maps. At the extremal tuple its evaluation $F$ is consequently
a finite sum of smooth cusp forms at this principal level. It
is bounded by RPI11. For each multiplicity map restrict to the
lowest compact subspaces. Its group equivariance gives the compact
equations, and differentiation gives the lowering equations: one
differentiates the translated smooth input vector and passes the
derivative through point evaluation by the map's separate Sobolev
continuity. The lowering operator kills the extremal input.
Equivalently each upper Borel derivative at the extremal tuple is
zero off the diagonal and the specified weight times the evaluation
on the diagonal. HCU5, or HCU19 on the restricted map, gives its
holomorphic type. The finite sum has the same type. Recover
a bounded component section $P$ from $F$ as follows. An
orthonormal frame is represented, up to its arithmetic stabilizer,
by a dictionary point $h$. Take
$F(h)/(c_{a(h)}\mu(h))$ there. Rational and principal-level
invariance and the archimedean weight give the same value for
the other representatives after the specified Borel weight is
applied. Gram orthogonalization expresses any frame as a unitary
translate of the reference frame times an upper triangular matrix.
On the compact triangular stabilizer the extremal law gives
exactly its Borel character, so extension by the Borel weight
is well defined. The local holomorphic representative in HCU19,
with the determinant powers from G77D2 removed, agrees on the
unitary translates and extends by that same Borel weight to the
full frame domain. The weighted norm there equals the norm on
orthonormal representatives and is bounded by the bound for $F$.
Extend $P$ by zero outside the domains. Its adelic reading is $F$.

We spell out the scope of the theta spanning argument. Let $P$
be any tuple of bounded holomorphic sections on these components
with their native weight laws and zero extension. Initially it
need not be the reading of an algebraic section. The period and
model clauses of $w.\mathrm{laws}$ at this $\alpha$ give the common
Gram matrix and polarization scalar on the components, the
principal congruence subgroup at $M$, and trivial arithmetic
weight characters. The integral symplectic period identification
on each component gives the self-dual trace lattice with its
marked half. Indeed its perfect alternating trace form identifies
the lattice with the trace dual as in PCL5--PCL6. In C76P3 use
$\alpha(\zeta_M)=\exp(2\pi i t/M)$; $t$ is still invertible
modulo $M$. The dual marking, generating tuple and separation of
component orbits in C76P3--C76P5 use just this congruence and
these actual period markings. Relabel the weight indices by the
type at $\alpha$ in C76P6. The stabilizer calculation C76P7
uses the same marked ideal and principal congruence.

Now use the analytical part of B76S2, taking $P_k$ in place
of $\mathrm{rd}(f)_k$: divide each by its component constant,
transport it to the representative marked lattice, extend to
its entire orbit by the vacuum multiplier, and set the other
orbits to zero. Disjointness and the stabilizer law give one
bounded holomorphic family with the prescribed Borel weight.
Since $n\ge4[E:\mathbb Q]r+4$, the reproducing formula I75K3
with B76S1 and the finite system L76F6 expands this family on
all marked lattices at once. Restrict to these representatives.
The same integral harmonic polynomial expansion, identity-cover
short sections of A76T2 and comparison S76G4 used in B76S2
produce
\[
        P = \sum_b d_b\,\mathrm{rd}_\alpha(f_b),\qquad d_b\in\mathbb C .
\]
The coefficient at a fixed generator is independent of the
component because the initial division by its component constant
and the final comparison cancel it. The equality off the
domains uses the zero extensions. G77D5 gives $f_b\in C_i$.
Here the conclusion is precisely a complex relation of component
functions; the last coefficient-descent step of B76S2 is unnecessary.
In particular this argument applies to the $P$ recovered from $F$
without first identifying its adelic transvection integrals.

Reading this relation adelically and dividing its coefficients by
the common nonzero Lie factor expresses $F$ as a complex sum of
$\phi_C(f_b)$. The core space law identifies the model Hermitian
space with the geometric one at every good restriction. G77D2
therefore gives $\mathrm{ev}_{v_0}j_{\rm Top}(f_b)=\phi_C(f_b)$.
Injectivity yields $y\in\operatorname{span}_{\mathbb C}j_{\rm Top}(C_i)$.
For independence, take any finite $K$-independent list in $C_i$.
Its component readings are independent after extension to
$\mathbb C$ by ICL8 on the good uniformization. A linear relation
between adelic readings vanishes on all dictionary frames, and
hence on every frame: the dictionary reaches each domain frame
up to an arithmetic mover and a triangular motion of nonzero
weight. The reading is zero outside the domain. The Lie factor,
component constants and $\mu$ are all nonzero. Thus the
$j_{\rm Top}$ images are also independent.

Finally $C_i$ has finite dimension, by its tied embedding in the
finite coordinate space. At a nongood restriction the chosen
conjugate top reading in G77D3 takes, at its one unit tuple, the
complex conjugates of the good readings. It therefore carries any
$K$-independent list to a complex-independent list as well.
The good space has dimension $\dim_K C_i$ by the preceding two
paragraphs; G77D3 gives the same complex dimension at its
conjugate. The images of a $K$-basis thus span at the nongood
restriction. This proves both clauses at every embedding.
% END G77D6

% BEGIN G77D7
\paragraph{G77D7. All seven fields of the supplied geometric half.}
Let $z$ be an original $R_3.\mathrm{Dat}$ at a face with its
pins, or simply take the underlying $B:\mathrm{GeoHalf}$ with
the principal-level and model laws and its base-row pins.
Suppose $8D.R.f\,D.R.r+4\le D.R.nV$.
Then $\mathrm{InGeo}(B)$ holds. On the original record the
weak-characterization level and model laws are the ones needed
in G77D2. Its two marked primes agree in order by the lying-over
argument of T76R1, using only the level law, the row pins and
the full weight/frame tie. In the source order the evaluation
and conjugate-evaluation fields are G77D2 and G77D3, and the
unit-vector norm argument of G77D2 gives $\mathrm{EvInj}$.
The two representation fields are G77D6 and G77D4. M76L2
applied to the core and weight laws gives $\mathrm{LatFull}$,
and $B.\mathrm{fw.dimMatch}$ identifies the stored row dimensions.
All seven statements concern this $B$ and its actual realization.

In particular the further $\mathrm{InGeo}$ premise in C76U1--C76U3
is now available on the original pinned datum at this rank; those
claims give $R.\mathrm{LUnif}$ with the indicated group laws.
G77D1 supplies $R.\mathrm{LIntegral}$ already on the underlying
geometric half. The estimates used here are a separate bound for
each fixed form from RPI11 and weighted estimates for the specified
short theta generators. The universal quotient-height statement
$\mathrm{CuspRapidDecay}$ has its own quantifiers.
% END G77D7
